\section{What causal reasoning benchmarks establish}
\label{gap29:frontier}
A banking analyst asks several different questions of an assistant. Can it
translate a request for higher deposits into an intervention and a measurable
target? Can it distinguish information known before an offer from information
created by the offer? Can it explain which arrows in a diagram are assumptions,
and recognise when several diagrams fit the same observations? A correct
answer to a narrative puzzle is useful evidence about some of these abilities,
but it does not establish all of them.

\subsection{Explanations deserve their own test}
Wang and Shen examine three constructed scenarios: coin operations, interrupted
file downloads and a hypothetical vaccination calculation. Four variants of
each scenario and three prompting approaches give 36 prompt conditions.
Their scoring distinguishes a correct result and explanation, a flawed
explanation, and an incorrect result. This is valuable because a correct
number can accompany a mistaken economic argument. The study reports mixed
changes from examples and guided reasoning, rather than a uniformly helpful
prompting intervention. \citep{gap29wangshen}

The terms used in a benchmark also need scrutiny. The paper calls reasoning
from an observed effect to a cause \emph{actual causality}. That task description
does not, by itself, specify a formal criterion for an individual's actual
cause. Its smallpox case asks which source produces more deaths in a stipulated
population. With the given numbers, vaccination reactions produce
$990{,}000(0.01)(0.01)=99$ deaths, while disease among unvaccinated people
produces $10{,}000(0.02)(0.20)=40$. Those population counts differ from the
per-person risks, $0.0001$ and $0.004$. Confusing a total with a risk is an
estimand error; the arithmetic alone does not establish a general
Simpson-reversal mechanism. These are deductions from the paper's stated
example, not new medical evidence.

The banking analogue is direct. A large customer group can contribute more
defaults while offering a smaller default risk per account. A deposit offer
can increase gross balances without increasing net contribution. An assistant
should state the denominator and economic target before presenting the
answer. The paper's prompts provide useful stress cases for that distinction.
The small scenario set and model configurations do not support a current
ranking of assistants for banking work. We have inspected the published
methods and results; we have not reproduced its model calls or inspected every
supplementary response.

\subsection{Scale does not remove uncertainty about the answer key}
MDC-Bench broadens the comparison to 8,422 questions across seven disciplines,
including business, economics and management. Its tasks distinguish reasoning,
intervention, decision and counterfactual questions. Questions are constructed
from causal narratives extracted from preprints, with model-generated
alternatives and a manual sampling check. The published protocol identifies
model versions, prompt templates, temperature and output limits.
\citep[Sections 2--3 and Appendix A]{gap29mdc}

Its harder condition requires two ranked answers instead of one. For example,
the reported Claude thinking configuration has average accuracy 84.56\% in
the basic condition and 59.25\% in the harder condition. These are different
response requirements on the same benchmark; their difference is not a
measured deterioration in one fixed banking task. The benchmark is useful
for comparing errors by causal structure and question type, provided its
answer-generation assumptions remain visible.

The detailed prompts reveal an important distinction. Reverse reasoning may
use necessity and sufficiency, or otherwise select an abductive explanation
judged probable from background knowledge. Intervention questions construct
an altered causal narrative; the answer alternatives are ranked by their
judged plausibility. These procedures do not supply a known structural
equation system, a randomised treatment effect or an identified posterior
over causal diagrams. Appendix A's sampling review helps detect poor
questions but does not turn model-generated causal statements into independent
economic truth. The authors also discuss weak long-chain relations and
distractors that may be too easy. \citep[Limitations; Tables 28--29 and
39--40]{gap29mdc}

\subsection{An agent architecture still needs an identified effect}
Chen combines language models, diffusion models and interacting agents in a
study of climate-policy sentiment. Section 6 describes 43,943 tweets and
comparisons with topic, language and diffusion baselines. It reports F1,
AUC, causal-chain precision, counterfactual error and simulated intervention
outcomes. Table 4 removes individual components to examine their contribution.
These are reported results of the paper; the model was not reproduced
locally. \citep[Sections 5.6 and 6; Tables 1--4]{gap29chen}

The causal interpretation needs more than those comparisons. Equation (16)
factorises a joint distribution according to a graph. Equation (17) defines
a conditional treatment contrast. Neither statement alone identifies that
contrast from the observed tweets. Equation (19) compares predicted and
reference counterfactual outcomes, but its accompanying text allows the
reference to come from annotations or simulated baselines. Section 6.1 does
not supply enough detail to reconstruct an independently identified
counterfactual for each tweet. F1 and AUC assess labelled prediction; they
do not repair that missing identification argument.

The useful proposal lesson is to separate an architecture from the evidence
for a decision. A bank may test whether a language model helps extract a
policy variable or whether a simulator represents a declared mechanism.
It still needs a defensible treatment assignment, observation model and
counterfactual comparison before attributing a financial gain to the policy.
The reported architecture is a research lead. Its stated improvements are
not evidence that the proposed banking service already delivers causal or
commercial benefits.

\subsection{A useful banking comparison}
The proposal should use such benchmarks to diagnose formulation and reasoning
failures, then add tasks whose consequences can be checked independently.
A compact questionnaire and an executable calculation are necessary
comparators. For a deposit question, all methods should receive the same
decision date, eligibility rule, alternative offer and contribution definition.
For servicing, all should see the same observed history and reasons for
dropout. For a diagram, an admissible set should remain admissible when the
observations cannot orient an edge.

The resulting assessment has distinct outcomes: whether the target is correct,
whether a proposed adjustment is justified, whether an unknown effect remains
unknown, whether the numerical calculation is correct, and how much analyst
time the assistance saves. Independent adjudication is needed for open-ended
economic formulations. Generated cases with known mechanisms can test
specific errors before an expensive bank exercise. The public experimental,
servicing and linkage studies in this volume provide numerical cases of that
kind, but no live assistant has yet been evaluated on them.

The inspected papers therefore sharpen the proposed experiment rather
than establishing a benefit from the proposed service. Repeated agreement
among prompts is especially weak evidence when their errors share a source,
as the repeated-measurement calculation shows. The useful next comparison
measures consequential mistakes and time on held-out tasks, with human
adjudication and a fixed cost allowance.
